\section{Introduction}

A sparse autoencoder (SAE) trained to reproduce the decoder directions of a parent SAE with fewer latents (meta-latents) is called a meta-SAE, which was introduced by \citet{bussmann2024metasae}. Meta-latents can be interpreted as representing more general concepts related to the latent's concept. For example, an ``Einstein'' latent might deconstruct into meta-latents like ``scientist'', ``prominent figures'', and ``Germany''. This idea was used by \citet{leask2025canonical} to argue that larger SAEs' latents are not always atomic, while SAE stitching showed that smaller SAEs' latents are incomplete in representing all features of the LLM. The underlying mechanism of meta-SAEs has been used to score the atomicity of SAEs \citep{korznikov2025ortsae}, to train atomicity-conscious SAEs \citep{levinson2026metasaes}, and to create cascaded SAEs for jointly learning a two-level dictionary \citep{zhao2026saepp}. However, all these studies assume reproducibility across meta-SAE seeds.

\citet{paulo2025sae} found partial seed reproducibility for ordinary SAEs, but no known studies to date investigate similarly for meta-SAEs. Studies have also typically not swept meta-SAE width and instead used conventional choices such as 1/4 \citep{bussmann2025matryoshka,korznikov2025ortsae} or 1/2 \citep{zhao2026saepp}. Thus, it matters whether the configuration adapted from \citet{leask2025canonical}---our reference---is reproducible across seeds to begin with and how seed-reproducible meta-SAEs are for various widths and underlying LLM sizes.

\begin{figure}[t]
  \centering
  \begin{tikzpicture}[
  font=\footnotesize,
  >={Stealth[length=3.4pt,width=3pt]},
  box/.style={draw, rounded corners=1.5pt, line width=0.5pt, align=center, inner xsep=3.5pt, inner ysep=2pt, minimum height=0.64in},
  real/.style={box, draw=acc!85!black, fill=acc!10},
  twin/.style={box, draw=black!80, fill=white},
  stack/.style={draw, rounded corners=1.5pt, line width=0.5pt, fill=white},
  flow/.style={->, line width=0.5pt},
  lane/.style={anchor=south west, inner sep=0pt, yshift=2.5pt, font=\footnotesize\bfseries},
]
\definecolor{acc}{RGB}{26,95,160}
\def\dy{1.0in}
\def\g{0.125in}
% real lane
\node[real, minimum width=0.56in] (r1) {Parent\\decoder\\rows\\$n\times d$};
\node[real, minimum width=0.50in, right=\g of r1] (r2) {Unit\\rows};
\node[real, minimum width=0.58in, right=\g of r2] (r3) {Train, val,\\test rows\\80/10/10};
\node[real, minimum width=0.84in, right=\g of r3] (r4) {Meta-SAE\\BatchTopK, $k{=}4$\\$wn$ meta-latents\\on train rows};
\node[real, minimum width=0.70in, right=\g of r4] (r5) {One-to-one\\matching\\28 seed pairs};
\node[real, minimum width=0.60in, right=\g of r5] (r6) {Real\\agreement};
% twin lane
\node[twin, minimum width=0.56in, below=\dy of r1.center, anchor=center] (t1) {Gaussian\\rows\\$n\times d$};
\node[twin, minimum width=0.50in, right=\g of t1] (t2) {Unit\\rows};
\node[twin, minimum width=0.58in, right=\g of t2] (t3) {Train\\rows\\80\%};
\node[twin, minimum width=0.84in, right=\g of t3] (t4) {Meta-SAE\\BatchTopK, $k{=}4$\\$wn$ meta-latents\\on train rows};
\node[twin, minimum width=0.70in, right=\g of t4] (t5) {One-to-one\\matching\\28 seed pairs};
\node[twin, minimum width=0.60in, right=\g of t5] (t6) {Twin\\agreement};
\begin{pgfonlayer}{background}
  % seeds as a stack of cards behind each meta-SAE box
  \foreach \n/\c in {r4/acc!85!black,t4/black!80} {
    \foreach \s in {3,6} {
      \draw[stack, draw=\c] ($(\n.north west)+(\s*0.7pt,\s pt)$) rectangle ($(\n.north east)+(-\s*0.7pt,\s pt-8pt)$);
    }
  }
\end{pgfonlayer}
\node[anchor=south, inner sep=1pt] at ([yshift=8pt]r4.north) {8 seeds};
\node[anchor=south, inner sep=1pt] at ([yshift=8pt]t4.north) {8 seeds};
\foreach \a/\b in {r1/r2,r2/r3,r3/r4,r4/r5,r5/r6,t1/t2,t2/t3,t3/t4,t4/t5,t5/t6}
  \draw[flow] (\a.east) -- (\b.west);
% twin rows take their mean and covariance from the parent's unit rows
\coordinate (mid) at ($(r2.south)!0.5!(t1.north)$);
\coordinate (cm) at (r2.south |- mid);
\coordinate (land) at ([xshift=-0.1in]t1.north east);
\draw[flow, dashed] (r2.south) -- (cm) -- (land |- mid) -- (land);
\node[anchor=west, inner sep=1pt] at ([xshift=3pt]cm) {mean, covariance};
% columns 1 to 3 of both lanes are shared by all widths and seeds
\coordinate (bl) at ([yshift=-6pt]t1.south west);
\coordinate (br) at ([yshift=-6pt]t3.south east);
\draw[line width=0.5pt, black!80] (bl |- t1.south) -- (bl) -- (br) -- (br |- t3.south);
\node[below, inner sep=1.5pt, text=black!80] at ($(bl)!0.5!(br)$) {shared by all widths and seeds};
% excess
\node[draw=acc!85!black, fill=acc, text=white, rounded corners=1.5pt, line width=0.5pt, align=center, inner xsep=3.5pt, inner ysep=2pt,
      minimum width=0.62in, minimum height=1.64in, anchor=west]
  (ex) at ($(r6.east)!0.5!(t6.east)+(\g,0)$) {Excess\\agreement\\[3pt]real${}-{}$twin};
\draw[flow] (r6.east) -- (r6.east -| ex.west);
\draw[flow] (t6.east) -- (t6.east -| ex.west);
\node[lane] at (r1.north west) {Real};
\node[lane] at (t1.north west) {Twin};
\end{tikzpicture}

  \caption{The excess agreement measurement pipeline for one parent SAE and meta-SAE width. The parent's $n$ decoder rows (dimension $d$) are scaled to unit length and, for a Gemma Scope 2 parent, split into training, validation, and test rows. Meta-SAEs with $wn$ meta-latents ($w$ is the width, as 1/4 to 1/128 for the main widths) train on the training rows with eight seeds that differ only in initial weights and shuffle order. Twin runs (lower lane) train on Gaussian rows but are validated and tested on the real rows.}
  \label{fig:pipeline}
\end{figure}

We run eight seeds for each meta-SAE configuration (width and LLM size) and then measure the agreement amongst them, where agreement is the mean cosine of Hungarian-matched pairs of decoder directions between pairs of unique seeds (Figure~\ref{fig:pipeline}). The seed-reproducibility of the reference configuration uses their parent SAE of GPT-2 small and meta-SAE settings adapted from theirs. The width-and-size sweep is done with five Gemma Scope 2 parent SAEs \citep{mcdougall2025gemmascope2} of the middle layers of respective Gemma 3 model sizes \citep[270M, 1B, 4B, 12B, and 27B;][]{gemmateam2025gemma3,googledeepmind2025gemma3card}, across six widths (1/4 to 1/128). To distinguish agreement from Hungarian-matched noise, we train noise twin meta-SAEs with the same settings as each meta-SAE configuration by training on unit-length Gaussian draws of the parent SAE's unit-length decoder directions and then record excess agreement (setup agreement minus noise twin agreement).

To our knowledge, this is the first width sweep of meta-SAE seeds' excess agreement. Our findings and contributions are as follows:
\begin{itemize}
  \item \textbf{Reference.} The reference configuration averages 64.0\% of meta-latents shared per seed pair (using one-to-one matching with at least 0.7 decoder cosine, which we call the recurring fraction) against 0.07\% for the noise twins. For the width-size sweep---whose parent SAEs are a different type, number of latents per hidden dimension, and underlying LLM, among other settings---the recurring fraction is less than a third.
  \item \textbf{Width by size.} For Gemma 3 270M and 1B, meta-SAE width 1/4 yields the highest excess agreement---although under 8\% of meta-latents recur---while smaller widths show a trend of lower, even non-positive, excess agreements. For 4B to 27B, widths 1/32 to 1/128 yield higher excess agreements than that of 1/4. Each main sweep width at every Gemma 3 model size has a meaningfully higher recurring fraction than its noise twins.
  \item \textbf{Noise twins.} While width 1/128 dominates raw agreement for four of the five Gemma 3 sizes, noise twin agreement over width is generally U-shaped for each model size. This heavily impacts which width has the best excess agreement when moving from 270M to 12B.
  \item \textbf{Controls.} Given equal training steps, comparisons between meta-SAEs trained on three randomly drawn quarters of the 4B parent's rows and on the released 16,384-latent Gemma Scope 2 SAE (same 4B layer) indicate that none of width, meta-latent count or row count alone determines excess agreement.
  \item \textbf{In-practice consequences.} When using meta-SAEs on different parent SAEs, one should consider excess agreement over noise twins instead of raw agreement and consider narrower widths (at least on Gemma Scope 2) for larger LLMs.
\end{itemize}
